\section{从鄗城到雒阳：复汉的名号与可见的边界}

建武元年六月，刘秀在鄗即皇帝位，改元建武。关中更始政权已在失序之中，河北则有另一套军政力量正在征发人力、安置官员；刘氏名号和新的年号由此成为将这些力量置于同一朝廷名义之下的办法。即位并没有使战争停住。它首先改变的是各地武装和官吏面对的名义：争夺不再只以一支军队或一个地方首领的名义展开，也以能否承认、供给或拒绝这个重新举起的汉廷展开。

同在这一年，刘秀定都雒阳。长安附近的战事和供给困境尚未解除，雒阳却更靠近其军政支持所系的关东与河北道路；都城的选择因而也是调度粮食、文书和军队的选择。汉魏洛阳城遗址所保存的城址、道路与宫城区域，能让后人看见这座城市在长时段中的政治重量，却不能把后来的遗存直接还原成建武元年某一座宫殿的形制，更不能据此断定当年营建的每一步。城址证明的是空间条件，不是仪式细节的替身。

\subsection{长安的终局，并非天下的终局}

建武二年，赤眉军进入长安，更始帝刘玄败亡。一个刘氏竞争中心由此消失，但关中并未立刻归入雒阳的命令之下。更始政权对将领、豪强和供给的约束已经破裂，赤眉军所面对的也不是一座可以长久供养大军的都城。长安的易手把战后人口、粮食和道路的破碎暴露出来；名义上的皇帝更替，并不能自动恢复地方秩序。

次年，赤眉主力在崤底一带向东汉军受降，作为全国性竞争力量的赤眉由此瓦解。传世记载能够相互印证这一过程及其大致年月，但战场的精确位置、受降人数和许多处置的细节并不一致。删去这些不能互证的数字，仍可看见一项实际变化：关东通向关中的道路已较能被东汉方面掌握，流民武装难以再依靠长安维持广泛的军事存在。雒阳朝廷于是得以把兵力和任命转向尚未归服的陇右与巴蜀。

\subsection{巴蜀与陇右：道路还没有合拢}

建武六年，公孙述在成都称帝，国号成家。巴蜀的山川、粮赋与地方网络，使成都能够成为一个不接受关东任命的政治中心；东汉的恢复因此不是从长安、雒阳之间的一条线向四方平推，而要面对西南的持久分立。成家的官制以及它在不同地区能够实际控制到何处，材料并不给出同样清楚的答案。可以确定的是，成都的称帝让巴蜀的归属成为重建赋役、交通与郡县任命时无法绕开的事情。

建武十年，隗嚣死后，其子隗纯降汉，陇右主要的割据力量被解除。隗氏一度在东汉与成家之间周旋，继承集团却未能延续原有的联盟。陇右并非地图上的一片空白：它连着关中、河西和通往西南的路径，控制这里意味着东汉可以在更广的道路网络上运兵、传令和筹给。降后处置的细节有异文，但这一通路的转折较为清楚；成家也由此少了一个可供牵连的东部力量。

建武十一年，朝廷下诏限制把战乱中掠卖的人继续充作奴婢，并令核实后释放。诏令所面对的是战争留下的逃户、掠卖与身份混乱；它重申了自由民身份的边界，也与重新编户、恢复赋役和治安相连。不过，法令的存在不能换算成已经获释的人数。地方豪强如何应对，官府核验到了何种程度，现有材料都不足以给出全国性的结论。

建武十二年，吴汉等军攻入成都，公孙述战死，成家灭亡。巴蜀于是被纳入东汉的郡县任命与赋役网络。攻城的经过和伤亡数字在叙事中尤其容易被铺陈，不能据此描画一幅毫无空隙的战场图；较稳的骨架是陇右平定之后，东汉能够集中力量入蜀，而长期孤立的成家终于失去其政权中心。军事上的压制完成后，恢复并未变成一张整齐的户籍册。

建武十五年的度田，正显出这一点。朝廷命郡国清查垦田、户口，地方的隐匿和官吏的争议随即进入记载。度田不是单纯的数字收集：若田地、人口和身份不能登记，赋役、治安和官府的日常运转便没有可据的基础。可是诏令与争议只能说明中央试图重新取得登记，不能证明清查覆盖了所有郡国，更不能让后人得到一项可与别朝直接比较的全国土地总数。东汉的重建，到这里仍是一种在道路、城邑和簿籍之间反复校定的过程。

\subsection{史料留下的亮处与阴影}

《后汉书》的帝纪和列传保存了即位、建都、关中战事、陇右归服、诏令与度田的主干；它使这些分散的行动能够连成朝廷的年序，却也是刘宋时范晔据较早史料写成的后出叙事，不能把它的政治评价或一切细部当作事件现场。《后汉纪》将东汉史事编为连续年月，可用来对读日期和异文；它在东晋编成，依托的材料链与《后汉书》并非彼此完全独立，二者相合也不等于每一个数字都已得到旁证。北宋《资治通鉴》可帮助观察前后相承的年序，同样不能替代更早材料。

雒阳的城址发掘给这段叙事增加了不同于史书的证据：都城并非抽象的名号，而有道路、城垣和宫城区域可以检验。然而发掘范围、断代与复原程度限制了它的发言；不能由一段遗址断言建武初年的城市规模，更不能由都城遗存推及全国社会。写到公孙述时，《华阳国志》所保留的巴蜀区域叙事尤其重要，它让成都不只作为雒阳史书中的远方城名出现；但这部后出的地方叙事也带着区域记忆和传本层次，不能单独裁定成家官制、疆域或攻城细节。本节因此只保留能够由这些材料交叉支撑的即位、迁都、政权兴亡、道路转折与诏令争议，不补写无人能证的言辞、仪式和人物心中唯一的动机。
